\chapter{排序、选择和区间问题}

排序不是为了把数组排好看，而是为了制造结构。无序输入里，两个元素之间没有稳定关系；排序之后，大小关系变成单调关系，重复元素挨在一起，区间端点可以按起点或终点扫描，双指针和二分才有施展空间。面试中很多题的关键不是“调用 \code{std::sort}”，而是说明排序之后哪个约束变简单了，排序代价是否值得，以及排序是否会破坏题目要求的原始下标或稳定性。

本章先讲排序的结构收益和比较器规则，再讲原地分类、区间合并、插入区间、区间贪心、快速选择和归并统计。区间题卡现在接在本章末尾，和排序选择讲义放在一起；贪心题中依赖区间终点的部分会在这里和堆贪心章互相呼应。

C++ 的 \code{std::sort} 通常是内省排序，综合快速排序、堆排序和插入排序思想，平均和最坏复杂度都能控制在 \complexity{O(n\log n)} 级别。刷题时可以放心使用它，但要理解比较器必须满足严格弱序。比较器不能写成 \code{a <= b}，也不能依赖会变化的外部状态，否则排序行为未定义。对结构体排序时，比较器应该只表达排序键，例如先按起点升序，再按终点升序。

\begin{lstlisting}[language=C++,caption={区间排序比较器：严格弱序}]
struct Interval {
    int start = 0;
    int end = 0;
};

void sort_intervals(std::vector<Interval>& intervals)
{
    std::sort(intervals.begin(), intervals.end(),
              [](const Interval& lhs, const Interval& rhs) {
                  if (lhs.start != rhs.start) {
                      return lhs.start < rhs.start;
                  }
                  return lhs.end < rhs.end;
              });
}
\end{lstlisting}

这个比较器在 \code{lhs.start == rhs.start && lhs.end == rhs.end} 时返回 \code{false}，满足严格弱序要求。很多线上题数据小，错误比较器也可能侥幸通过，但这是 C++ 基本功问题。排序后，如果题目仍需要原始下标，就不能只排序值；要把值和下标一起保存成 \code{pair} 或结构体。

颜色分类是荷兰国旗问题。题目只包含 0、1、2 三种值，要求原地排序。直接调用排序当然能过，复杂度 \complexity{O(n\log n)}；计数排序也能 \complexity{O(n)}，但需要两次扫描。更经典的一次扫描做法是维护三个区域：左侧全是 0，中间是已经确认的 1，右侧全是 2，未知区域在中间游动。指针 \code{low} 指向下一个 0 应放的位置，\code{mid} 扫描未知区域，\code{high} 指向下一个 2 应放的位置。

\begin{lstlisting}[language=C++,caption={LeetCode 75 颜色分类：三指针维护四段区域}]
void sort_colors(std::vector<int>& nums)
{
    constexpr int RED = 0;
    constexpr int WHITE = 1;
    constexpr int BLUE = 2;

    int low = 0;
    int mid = 0;
    int high = static_cast<int>(nums.size()) - 1;

    while (mid <= high) {
        if (nums[mid] == RED) {
            std::swap(nums[low], nums[mid]);
            ++low;
            ++mid;
        } else if (nums[mid] == WHITE) {
            ++mid;
        } else if (nums[mid] == BLUE) {
            std::swap(nums[mid], nums[high]);
            --high;
        }
    }
}
\end{lstlisting}

这题最容易错在交换 2 之后是否移动 \code{mid}。答案是不能移动，因为从 \code{high} 换过来的元素还没有检查，可能是 0、1 或 2。交换 0 后可以移动 \code{mid}，因为 \code{low} 左边已经处理完，换到 \code{mid} 的元素来自已确认的 1 区或当前位置本身。这个差别正是不变量的体现：\code{[0, low)} 是 0，\code{[low, mid)} 是 1，\code{(high, n)} 是 2，\code{[mid, high]} 未知。

合并区间展示排序如何把二维关系变成线性扫描。暴力方法是反复找任意两个相交区间合并，直到不能合并为止；每次合并后都可能重新扫描全部区间，复杂度很高。按起点排序后，相交关系只需要和当前结果的最后一个区间比较。因为后续区间起点只会越来越大，如果它不能和最后区间相交，就不可能和更早已经结束的区间相交。

\begin{lstlisting}[language=C++,caption={LeetCode 56 合并区间：排序后只看结果尾部}]
std::vector<std::vector<int>> merge_intervals(
    std::vector<std::vector<int>> intervals)
{
    if (intervals.empty()) {
        return {};
    }

    std::sort(intervals.begin(), intervals.end(),
              [](const std::vector<int>& lhs, const std::vector<int>& rhs) {
                  return lhs[0] < rhs[0];
              });

    std::vector<std::vector<int>> merged;
    merged.push_back(intervals.front());

    for (int i = 1; i < static_cast<int>(intervals.size()); ++i) {
        std::vector<int>& last = merged.back();
        const std::vector<int>& current = intervals[i];
        if (current[0] <= last[1]) {
            last[1] = std::max(last[1], current[1]);
        } else {
            merged.push_back(current);
        }
    }

    return merged;
}
\end{lstlisting}

这里参数按值传入，是因为函数内部排序会改变区间顺序。如果调用者不需要保留原顺序，传引用也可以。复杂度是排序 \complexity{O(n\log n)} 加扫描 \complexity{O(n)}。易错点是相交条件：如果区间是闭区间，\code{current[0] <= last[1]} 表示相交或接触；如果题目定义半开区间，就要重新审视等号是否成立。区间题一定要先明确端点语义。

插入区间与合并区间类似，但输入已经有序且不重叠。暴力方法可以把新区间放进去重新排序再合并，复杂度仍可接受，但没有利用已有条件。线性做法分三段：先加入所有结束在新区间之前的区间；再把所有和新区间相交的区间合并进新区间；最后加入剩余区间。这种写法体现了区间扫描的分段思维。

\begin{lstlisting}[language=C++,caption={LeetCode 57 插入区间：前段、中段、后段}]
std::vector<std::vector<int>> insert_interval(
    const std::vector<std::vector<int>>& intervals,
    std::vector<int> new_interval)
{
    std::vector<std::vector<int>> answer;
    int index = 0;
    const int n = static_cast<int>(intervals.size());

    while (index < n && intervals[index][1] < new_interval[0]) {
        answer.push_back(intervals[index]);
        ++index;
    }

    while (index < n && intervals[index][0] <= new_interval[1]) {
        new_interval[0] = std::min(new_interval[0], intervals[index][0]);
        new_interval[1] = std::max(new_interval[1], intervals[index][1]);
        ++index;
    }
    answer.push_back(new_interval);

    while (index < n) {
        answer.push_back(intervals[index]);
        ++index;
    }

    return answer;
}
\end{lstlisting}

这题的优化条件是“原区间已经按起点排序且互不重叠”。如果这个条件不存在，就必须先排序或换模型。第一段的判断是 \code{intervals[index][1] < new_interval[0]}，说明当前区间完全在新区间左边；第二段的判断是 \code{intervals[index][0] <= new_interval[1]}，说明存在重叠。两个等号同样取决于闭区间语义。面试中把三段说清楚，代码通常很自然。

无重叠区间是区间贪心题。题目要求删除最少区间，使剩余区间互不重叠。等价地，可以保留最多数量的不重叠区间。暴力方法是枚举所有区间子集，检查是否互不重叠，复杂度指数级。贪心优化依据是：每次选择结束最早的区间，会给后面留下最大的空间。按结束时间排序，能选就选，不能选就跳过。

\begin{lstlisting}[language=C++,caption={LeetCode 435 无重叠区间：按结束时间贪心保留最多区间}]
int erase_overlap_intervals(std::vector<std::vector<int>> intervals)
{
    if (intervals.empty()) {
        return 0;
    }

    std::sort(intervals.begin(), intervals.end(),
              [](const std::vector<int>& lhs, const std::vector<int>& rhs) {
                  return lhs[1] < rhs[1];
              });

    int kept = 0;
    int current_end = std::numeric_limits<int>::min();
    for (const std::vector<int>& interval : intervals) {
        if (interval[0] >= current_end) {
            ++kept;
            current_end = interval[1];
        }
    }

    return static_cast<int>(intervals.size()) - kept;
}
\end{lstlisting}

这个贪心的交换论证是：假设某个最优解第一步选择了一个结束更晚的区间，而存在另一个兼容区间结束更早。把最优解里的第一个区间替换成结束更早的区间，不会减少后续可选空间，因此不会让答案变差。反复替换，就得到按结束时间贪心的解。区间贪心题不要只说“排序后贪心”，必须说明排序键为什么是结束时间而不是开始时间。

数组中的第 K 大元素可以排序解决，时间 \complexity{O(n\log n)}。如果只需要第 K 大，不需要完整有序，快速选择更合适。它借鉴快速排序的分区思想：选择一个枢轴，把大于枢轴的放左边，小于枢轴的放右边；如果枢轴最终位置正好是 \code{k - 1}，答案找到；如果位置偏左或偏右，只在一侧继续。平均复杂度 \complexity{O(n)}。

\begin{lstlisting}[language=C++,caption={LeetCode 215 数组中的第 K 个最大元素：迭代快速选择}]
int find_kth_largest_quickselect(std::vector<int> nums, int k)
{
    const int target = k - 1;
    int left = 0;
    int right = static_cast<int>(nums.size()) - 1;

    while (left <= right) {
        const int pivot = nums[right];
        int store = left;
        for (int i = left; i < right; ++i) {
            if (nums[i] > pivot) {
                std::swap(nums[store], nums[i]);
                ++store;
            }
        }
        std::swap(nums[store], nums[right]);

        if (store == target) {
            return nums[store];
        }
        if (store < target) {
            left = store + 1;
        } else {
            right = store - 1;
        }
    }

    return nums[target];
}
\end{lstlisting}

这里按值传参是刻意的，因为快速选择会重排数组。如果题目允许修改输入，可以传引用减少一次复制。代码使用最后一个元素作枢轴，最坏情况下会退化到 \complexity{O(n^2)}；工业实现通常会随机化枢轴或使用更稳健的选择算法。面试中要主动说出这个风险。快速选择的正确性来自分区后的位置含义：左侧元素都比枢轴大，因此枢轴排名已经确定，不需要管左右两侧内部是否有序。

排序和选择题经常和去重绑定。三数之和先排序，是为了让固定一个数后，剩余两数可以用双指针线性查找；合并区间先排序，是为了让所有可能相交的区间在扫描时相邻；第 K 大使用分区，是为了只确定目标排名而不完整排序。三者共同点是：排序或分区不是目的，而是让后续操作只看局部信息。

归并排序在刷题里经常不是为了排序本身，而是为了在“合并两个有序段”的过程中统计跨区间关系。逆序对就是典型例子。暴力方法枚举所有 \code{i < j}，检查 \code{nums[i] > nums[j]}，复杂度 \complexity{O(n^2)}。归并排序把数组分成左右两半，递归意义上左右内部逆序对已经统计完；合并时，只需要统计左半元素和右半元素形成的跨区间逆序对。由于左右两边已经有序，若 \code{left[i] > right[j]}，那么从 \code{i} 到左半末尾的所有元素都大于 \code{right[j]}。

\begin{lstlisting}[language=C++,caption={逆序对：归并过程中统计跨区间关系}]
std::int64_t count_reverse_pairs(std::vector<int> nums)
{
    std::vector<int> buffer(nums.size());
    std::int64_t answer = 0;

    for (int width = 1; width < static_cast<int>(nums.size()); width *= 2) {
        for (int left = 0; left < static_cast<int>(nums.size()); left += 2 * width) {
            const int mid = std::min(left + width, static_cast<int>(nums.size()));
            const int right = std::min(left + 2 * width, static_cast<int>(nums.size()));

            int i = left;
            int j = mid;
            int out = left;
            while (i < mid && j < right) {
                if (nums[i] <= nums[j]) {
                    buffer[out++] = nums[i++];
                } else {
                    answer += mid - i;
                    buffer[out++] = nums[j++];
                }
            }
            while (i < mid) {
                buffer[out++] = nums[i++];
            }
            while (j < right) {
                buffer[out++] = nums[j++];
            }
            for (int k = left; k < right; ++k) {
                nums[k] = buffer[k];
            }
        }
    }

    return answer;
}
\end{lstlisting}

这里使用自底向上的迭代归并，避免递归。 \code{width} 表示当前已经有序的段长度，每轮把相邻两段合并成更长有序段。统计逆序对的关键行是 \code{answer += mid - i}：当右侧当前元素更小，左侧从 \code{i} 到 \code{mid - 1} 的所有元素都会和它形成逆序对。这个推理依赖左半段已经有序。归并类统计题的复盘要问：跨区间关系能不能在两个有序段合并时一次性数出来？如果能，通常可以从平方复杂度降到 \complexity{O(n\log n)}。

会议室问题把区间排序和堆结合起来。会议室 I 问一个人能否参加所有会议，只要按开始时间排序，检查相邻会议是否重叠即可。会议室 II 问至少需要多少会议室，就不能只看相邻会议，因为多个会议可能同时进行。暴力方法对每个会议扫描当前所有房间，找能复用的房间；优化方法用小顶堆保存每个房间当前会议的结束时间。新会议开始时，如果最早结束的房间已经空出来，就复用它；否则需要新房间。

\begin{lstlisting}[language=C++,caption={LeetCode 253 会议室 II：小顶堆保存房间结束时间}]
int min_meeting_rooms(std::vector<std::vector<int>> intervals)
{
    if (intervals.empty()) {
        return 0;
    }

    std::sort(intervals.begin(), intervals.end(),
              [](const std::vector<int>& lhs, const std::vector<int>& rhs) {
                  return lhs[0] < rhs[0];
              });

    std::priority_queue<int, std::vector<int>, std::greater<int>> endings;
    for (const std::vector<int>& meeting : intervals) {
        if (!endings.empty() && endings.top() <= meeting[0]) {
            endings.pop();
        }
        endings.push(meeting[1]);
    }

    return static_cast<int>(endings.size());
}
\end{lstlisting}

堆顶表示当前所有占用房间中最早结束的会议。如果它都晚于新会议开始，那么其他房间结束更晚，也不能复用；如果它不晚于新会议开始，就弹出并把该房间的新结束时间压入。复杂度是排序 \complexity{O(n\log n)} 加每个会议一次堆操作。注意这里一次只弹出一个结束时间就够了，因为每个新会议只需要一个房间；如果题目问某个时间点同时进行会议数量，可能会弹出所有已结束会议。

差分数组是区间更新题的另一种排序替代思路。如果题目反复对区间加值，再询问每个位置最终值，暴力对每次更新逐个位置加，复杂度 \complexity{O(qn)}。差分数组保存相邻元素变化量：对区间 \code{[left, right]} 加 \code{value}，只需要 \code{diff[left] += value}，\code{diff[right + 1] -= value}。最后做一次前缀和还原所有位置。

\begin{lstlisting}[language=C++,caption={LeetCode 1109 航班预订统计：差分数组处理区间加法}]
std::vector<int> corp_flight_bookings(const std::vector<std::vector<int>>& bookings,
                                      int n)
{
    std::vector<int> diff(n + 1, 0);
    for (const std::vector<int>& booking : bookings) {
        const int left = booking[0] - 1;
        const int right = booking[1] - 1;
        const int seats = booking[2];
        diff[left] += seats;
        if (right + 1 < n) {
            diff[right + 1] -= seats;
        }
    }

    std::vector<int> answer(n, 0);
    int current = 0;
    for (int i = 0; i < n; ++i) {
        current += diff[i];
        answer[i] = current;
    }
    return answer;
}
\end{lstlisting}

差分的正确性来自“影响从 left 开始，在 right 之后结束”。它把区间内每个点都加的操作，变成两个边界事件。扫描线也有类似思想：把区间起点看成 +1 事件，终点看成 -1 事件，排序后扫描事件就能得到当前重叠数量。差分适合坐标连续且范围可开数组的场景；扫描线适合坐标很大或稀疏，需要排序事件的场景。

计数排序和桶排序利用值域信息。若数组长度很大但值域很小，比较排序的 \complexity{O(n\log n)} 下界不再是必须承受的，因为我们不靠两两比较区分所有排列，而是统计每个值出现次数。颜色分类可以看成值域只有 3 的计数排序变体；前 K 高频元素也可以在统计频次后按频次分桶。桶排序的风险是空间由值域或频次范围决定，不能在值域巨大时盲目开数组。

\begin{lstlisting}[language=C++,caption={LeetCode 347 前 K 高频元素：桶按频次收集}]
std::vector<int> top_k_frequent_bucket(const std::vector<int>& nums, int k)
{
    std::unordered_map<int, int> frequency;
    frequency.reserve(nums.size());
    for (const int x : nums) {
        ++frequency[x];
    }

    std::vector<std::vector<int>> buckets(nums.size() + 1);
    for (const auto& [value, count] : frequency) {
        buckets[count].push_back(value);
    }

    std::vector<int> answer;
    for (int count = static_cast<int>(buckets.size()) - 1;
         count >= 0 && static_cast<int>(answer.size()) < k;
         --count) {
        for (const int value : buckets[count]) {
            answer.push_back(value);
            if (static_cast<int>(answer.size()) == k) {
                break;
            }
        }
    }
    return answer;
}
\end{lstlisting}

这个版本的时间复杂度是 \complexity{O(n)}，空间复杂度也是 \complexity{O(n)}，因为最大频次不超过数组长度。它比堆更快吗？理论上是线性的，但常数和内存占用更高，且输出顺序不固定。实际面试可以先讲堆，因为堆适用范围更广；再补桶排序作为值域或频次范围明确时的优化。高质量答案不是只追求最低复杂度，而是解释方案适用条件。

相对排序数组展示“值域有限时不一定要比较排序”。题目给两个数组，第二个数组给出部分元素的相对顺序，要求第一个数组中这些元素按第二个数组顺序排列，其余元素升序排列。暴力做法可以写自定义比较器：若两个元素都在第二个数组中，就比较它们的排名；若只有一个在第二个数组中，它排前面；否则按值升序。这个做法清晰，但每次比较都要查询排名，复杂度仍是排序的 \complexity{O(n\log n)}。如果值域较小，可以计数后按规则输出，复杂度接近线性。

\begin{lstlisting}[language=C++,caption={LeetCode 1122 数组的相对排序：计数后按给定顺序输出}]
std::vector<int> relative_sort_array(const std::vector<int>& arr1,
                                     const std::vector<int>& arr2)
{
    constexpr int MAX_VALUE = 1000;
    std::array<int, MAX_VALUE + 1> count{};
    for (const int x : arr1) {
        ++count[x];
    }

    std::vector<int> answer;
    answer.reserve(arr1.size());
    for (const int x : arr2) {
        while (count[x] > 0) {
            answer.push_back(x);
            --count[x];
        }
    }
    for (int value = 0; value <= MAX_VALUE; ++value) {
        while (count[value] > 0) {
            answer.push_back(value);
            --count[value];
        }
    }
    return answer;
}
\end{lstlisting}

这段代码的前提是题目给定值域不超过 1000，所以 \code{MAX_VALUE} 是题目约束，不是随便拍脑袋。若值域很大，就不能开这么大的数组，应改用哈希表统计，再对剩余元素排序。这个题训练的是“比较排序下界的适用范围”：当元素值域和输出顺序被额外约束时，我们可以不通过比较确定所有相对顺序。

按奇偶排序、按自定义规则排序、把负数放前面等题，看起来简单，其实都在考“稳定性”。如果题目要求保持同类元素的原始相对顺序，就需要稳定划分，可以新开数组，或者用 \code{stable_partition}。如果不要求稳定，就可以用左右指针原地交换。原地不稳定划分的复杂度低、空间少，但会改变相对顺序；稳定划分保留顺序，但通常需要额外空间或更复杂的移动。面试中主动问清稳定性，是比直接写代码更成熟的表现。

\begin{lstlisting}[language=C++,caption={LeetCode 283 移动零：稳定划分保持相对顺序}]
std::vector<int> stable_partition_by_sign(const std::vector<int>& nums)
{
    std::vector<int> answer;
    answer.reserve(nums.size());
    for (const int x : nums) {
        if (x < 0) {
            answer.push_back(x);
        }
    }
    for (const int x : nums) {
        if (x >= 0) {
            answer.push_back(x);
        }
    }
    return answer;
}
\end{lstlisting}

这不是最省空间的写法，但语义非常清晰。若题目必须原地且稳定，问题会变难，因为把一个负数移动到前面可能要整体搬移中间元素。很多面试题默认不要求稳定，允许左右指针交换；但如果业务语义中原始顺序有意义，例如日志时间顺序、用户输入顺序，就不能随便破坏。排序和划分题一定要把“允许改变什么”说清楚。

数组中的逆序对除了归并计数，还可以用树状数组。归并排序从“跨左右两半的逆序对”出发；树状数组从“扫描到当前数时，之前有多少数比它大”出发。由于数值可能很大或为负，需要先离散化，把值映射到排名。扫描数组时，查询已经出现过的元素总数减去小于等于当前排名的数量，就是以当前元素为右端的逆序对数量；然后把当前排名出现次数加一。

\begin{lstlisting}[language=C++,caption={逆序对：离散化加树状数组}]
class FenwickTree {
public:
    explicit FenwickTree(int size) : tree_(size + 1, 0) {}

    void add(int index, int delta)
    {
        for (int i = index; i < static_cast<int>(this->tree_.size());
             i += i & -i) {
            this->tree_[i] += delta;
        }
    }

    int sum(int index) const
    {
        int result = 0;
        for (int i = index; i > 0; i -= i & -i) {
            result += this->tree_[i];
        }
        return result;
    }

private:
    std::vector<int> tree_;
};

std::int64_t reverse_pairs_by_fenwick(const std::vector<int>& nums)
{
    std::vector<int> values = nums;
    std::sort(values.begin(), values.end());
    values.erase(std::unique(values.begin(), values.end()), values.end());

    FenwickTree tree(static_cast<int>(values.size()));
    std::int64_t answer = 0;
    for (int i = 0; i < static_cast<int>(nums.size()); ++i) {
        const int rank = static_cast<int>(
            std::lower_bound(values.begin(), values.end(), nums[i]) -
            values.begin()) + 1;
        const int seen = i;
        const int not_greater = tree.sum(rank);
        answer += seen - not_greater;
        tree.add(rank, 1);
    }
    return answer;
}
\end{lstlisting}

树状数组适合动态维护前缀计数。它的下标从 1 开始，\code{i & -i} 表示当前节点覆盖区间的长度。这个结构比归并计数更抽象，但能迁移到很多“动态排名、区间频次、前缀和修改查询”问题。面试中如果只考逆序对，归并排序更容易解释；如果题目有多次更新或在线查询，树状数组更有优势。这里再次体现：排序不是唯一制造顺序的方法，离散化加索引树也能制造排名结构。

扫描线是区间题的重要模型。区间重叠数量、天际线、会议室、航班预订，都可以看成一系列事件按坐标发生。区间开始是一个加事件，区间结束是一个减事件。排序事件后从左到右扫描，当前累计值就是当前位置的活跃区间数量。对于闭区间和半开区间，开始和结束在同一坐标时的处理顺序不同；例如会议室通常认为一个会议在时间 t 结束，另一个会议在时间 t 开始可以复用同一房间，所以结束事件应该先于开始事件处理。

\begin{lstlisting}[language=C++,caption={会议室数量：扫描线事件排序}]
int min_meeting_rooms_by_events(const std::vector<std::vector<int>>& intervals)
{
    std::vector<std::pair<int, int>> events;
    events.reserve(intervals.size() * 2);
    for (const std::vector<int>& interval : intervals) {
        events.emplace_back(interval[0], 1);
        events.emplace_back(interval[1], -1);
    }

    std::sort(events.begin(), events.end(),
              [](const auto& lhs, const auto& rhs) {
                  if (lhs.first != rhs.first) {
                      return lhs.first < rhs.first;
                  }
                  return lhs.second < rhs.second;
              });

    int active = 0;
    int answer = 0;
    for (const auto [time, delta] : events) {
        (void)time;
        active += delta;
        answer = std::max(answer, active);
    }
    return answer;
}
\end{lstlisting}

这里 \code{-1} 事件在同一时间排在 \code{+1} 前面，所以会议结束先释放房间。若题目是闭区间覆盖点，结束和开始的同点处理可能要反过来。扫描线的难点不是代码，而是事件语义：一个事件表示影响从哪里开始，在哪里结束，边界是否包含。差分数组是坐标连续且范围可开数组时的扫描线；事件排序是坐标稀疏或很大时的扫描线。

用最少箭引爆气球是区间贪心。每个气球是一个水平区间，一支箭射在某个坐标可以引爆所有覆盖该坐标的气球。暴力想法是尝试每个可能射击位置，但坐标范围巨大。排序后贪心：按右端点升序排列，第一支箭射在第一个气球的右端点，这个选择最不容易影响后续，因为它尽可能靠右，同时仍能引爆当前气球。之后扫描气球，若当前气球起点大于当前箭位置，说明必须新射一箭。

\begin{lstlisting}[language=C++,caption={LeetCode 452 最少箭引爆气球：按右端点贪心}]
int find_min_arrow_shots(std::vector<std::vector<int>> points)
{
    if (points.empty()) {
        return 0;
    }
    std::sort(points.begin(), points.end(),
              [](const std::vector<int>& lhs, const std::vector<int>& rhs) {
                  return lhs[1] < rhs[1];
              });

    int arrows = 1;
    int arrow_position = points[0][1];
    for (int i = 1; i < static_cast<int>(points.size()); ++i) {
        if (points[i][0] <= arrow_position) {
            continue;
        }
        ++arrows;
        arrow_position = points[i][1];
    }
    return arrows;
}
\end{lstlisting}

正确性可以用交换论证说明：对于最早结束的气球，任何可行解都必须有一支箭射在它的区间内。把这支箭移动到该气球右端点，不会让它失去对已覆盖气球的覆盖能力，反而给后续区间留下更大机会。因此选择最早右端点是安全的。这个思路也适用于无重叠区间中“保留最多区间”：按结束时间选择最早结束的区间，给后面留下空间。

区间列表交集是双指针题。两个区间列表各自有序且互不重叠。暴力两两比较复杂度 \complexity{O(mn)}。双指针优化来自有序性：当前两个区间的交集左端是较大起点，右端是较小终点；若左端不大于右端，就有交集。然后谁的终点更早，谁就不可能再和对方后续区间相交，移动谁的指针。

\begin{lstlisting}[language=C++,caption={LeetCode 986 区间列表交集：移动较早结束的一方}]
std::vector<std::vector<int>> interval_intersection(
    const std::vector<std::vector<int>>& first,
    const std::vector<std::vector<int>>& second)
{
    std::vector<std::vector<int>> answer;
    int i = 0;
    int j = 0;
    while (i < static_cast<int>(first.size()) &&
           j < static_cast<int>(second.size())) {
        const int left = std::max(first[i][0], second[j][0]);
        const int right = std::min(first[i][1], second[j][1]);
        if (left <= right) {
            answer.push_back({left, right});
        }

        if (first[i][1] < second[j][1]) {
            ++i;
        } else {
            ++j;
        }
    }
    return answer;
}
\end{lstlisting}

这题的推理比代码重要：终点较早的区间在当前比较后已经没有机会和另一个当前区间或它后面的区间形成新的交集，因为后续区间起点只会更大。因此可以安全丢弃它。这个“丢弃谁”的逻辑和合并区间、插入区间、会议室复用都属于同一类端点单调推理。

日志重排和版本号比较这类题看起来不是传统算法，但同样考排序规则。日志重排要求字母日志排在数字日志前，字母日志按内容排序，内容相同按标识符排序，数字日志保持原始相对顺序。若直接写比较器处理数字日志之间返回 \code{false}，还需要使用稳定排序才能保持它们原顺序；或者把字母日志单独取出排序，再追加数字日志。这里的关键是理解 \code{std::sort} 不稳定，不能指望相等元素原顺序不变。

\begin{warningbox}
排序比较器有一个底线：不能让两个元素互相都“小于”等关系成立，也不能让相等元素比较返回真。比较器依赖外部可变状态、随机数、递增计数器，都会破坏排序算法假设。若题目要求保持相对顺序，请使用稳定排序或显式分组，不要假设普通排序“碰巧没改顺序”。
\end{warningbox}

区间和排序题的测试要覆盖边界语义。合并区间要测相邻是否合并、完全包含、乱序输入；会议室要测一个会议结束另一个会议开始；最少箭要测端点相同的大坐标；快速选择要测重复值、K 为 1、K 为数组长度；计数排序要测值域边界。很多排序题代码看起来一行 \code{sort} 很短，但真正的错误都藏在比较器和边界定义里。

\begin{labbox}
排序区间实验：第一，把会议室问题分别用堆和扫描线写一遍，用同一组边界用例比较答案。第二，把无重叠区间、最少箭、会议安排放在一起复盘，写出它们为什么都偏好按右端点排序。第三，给相对排序构造值域很大但元素很少的用例，比较计数数组和哈希统计的空间差异。第四，故意写一个 \code{a <= b} 的错误比较器，观察小数据可能不报错但语义已经未定义。
\end{labbox}

\begin{principlebox}{排序题复盘模板}
先写清楚排序键，再说明排序后获得了什么性质：相同值相邻、端点单调、排名确定、候选可淘汰，还是可以双指针。然后说明是否需要保留原下标，比较器是否满足严格弱序，排序是否改变输入。最后比较完整排序、堆、快速选择、计数排序等方案的适用范围。能说清这些，排序题才不是“先 sort 一下试试”。
\end{principlebox}

\section{排序、选择和分治的系统讲解}

排序题表面是在排列元素，实际是在制造顺序信息。没有顺序时，两数比较、区间合并、去重、前驱后继都可能需要线性寻找；排序以后，许多关系会变成相邻关系或单调关系。三数之和能够从三层枚举降到固定一数加双指针，是因为排序把剩余两数的和变成可单调调整的量。合并区间能够线性扫描，是因为按起点排序后，所有可能与当前结果尾部相交的区间都会连续出现。排序不是免费的，复杂度通常是 O(n log n)，所以必须说明排序换来了什么查询能力。

比较排序的下界来自信息论：n 个不同元素有 n! 种可能顺序，比较一次最多提供一位左右的信息，因此最坏至少需要数量级为 n log n 的比较。这个结论告诉我们，想突破 n log n，必须利用额外条件，例如值域有限、元素是整数、只需要第 K 个而不是完整顺序、或者允许随机化平均复杂度。计数排序利用小值域，桶排序利用分布，基数排序利用数字位，快速选择利用只关心目标排名。这些算法不是比标准排序更高级，而是使用了更强前提。

快速排序的核心是 partition：选择一个枢轴，把小于它的元素放一边，大于它的元素放另一边。partition 完成后，枢轴位置已经确定，左右两边再分别排序。快速选择只进入目标排名所在的一边，因此平均时间为线性。这里的平均性依赖枢轴足够随机；若每次都选到极端枢轴，最坏会退化到平方。面试中写快速选择必须主动说明随机化或三数取中，不然复杂度承诺不完整。

归并排序和快速排序的思维不同。归并排序先分再合，合并两个有序段时保证稳定性，最坏复杂度稳定为 O(n log n)，但需要额外空间。链表排序常用归并，因为链表合并不需要随机访问，节点重连成本低；数组排序常用 introsort 一类混合策略，在快速排序、堆排序、插入排序之间切换，兼顾平均性能和最坏情况。理解这些取舍后，才能解释为什么 C++ 的 std::sort 不承诺稳定，而 std::stable\_sort 需要更多空间。

选择第 K 大时有三种主线。完整排序最简单，适合数据量不大或后续还需要有序结果；大小为 K 的堆适合 K 远小于 n 或数据流持续到来；快速选择适合一次性数组、需要平均线性。三种方法没有绝对优劣，差别在输出需求、数据是否在线、实现风险和最坏复杂度。高质量题解应该比较它们，而不是只写一种。

排序题的边界往往出在比较器。比较器必须满足严格弱序，不能写出自相矛盾的关系。区间覆盖题中同起点要按终点降序，合并区间按起点升序即可，会议室按开始时间或结束时间有不同含义。比较器写错时，代码看似能跑，但排序结果不满足算法不变量。每次写比较器都要用三组数据验证：相同主键、完全相同元素、主键相反顺序。

实验建议：实现完整排序、大小为 K 的堆、快速选择三种第 K 大算法，用随机数组、几乎有序数组、逆序数组和大量重复数组测试。记录比较次数或运行时间，观察 K 很小时堆的优势、需要完整有序时排序的优势、随机化快速选择的波动。再把快速选择枢轴固定为首元素，用有序数组制造最坏情况，理解为什么工程库不会使用裸快速排序。

\begin{principlebox}{本节复盘}
阅读本节后，请回到相邻章节的 LeetCode 案例，例如 LeetCode 875 爱吃香蕉、LeetCode 72 编辑距离、LeetCode 743 网络延迟时间、LeetCode 215 数组第 K 大，把暴力瓶颈、优化依据、状态不变量、C++ 容器成本和边界测试重新写一遍。能从这些具体题迁移到变体题，才算真正掌握。
\end{principlebox}

排序和区间深度题卡 不再把所有题目堆成一个清单，而是按题型分成章内大节。每一节先讲分流规则，再列代表案例，最后逐题做单点分析。

\section{区间、排序与合并：题型分流与案例}
这一节先把同一题型的常见场景拆开讲清楚，再进入具体题目。学习顺序是：先写暴力基线，定位最贵的重复工作；再根据题目条件选择能维护不变量的数据结构；最后用多个案例比较边界、反例和变体。

\begin{principlebox}{图示：区间、排序与合并推导链}
\begin{tabular}{@{}p{0.18\linewidth}p{0.74\linewidth}@{}}
暴力基线 & 枚举候选、完整模拟或逐个检查，先保证覆盖全部可能答案。\\
瓶颈定位 & 瓶颈是没有利用排序后相邻关系。区间按起点或终点排序后，很多远离的区间不可能再影响当前决策。\\
结构选择 & 合并按起点排序，调度按终点排序，覆盖题常用起点升序且终点降序。扫描时维护当前最远右端、结果尾区间或已选区间终点。\\
不变量 & 结果数组中的区间已经互不重叠且有序；当前扫描区间只可能影响结果最后一个区间或当前边界。\\
代码落地 & 比较器必须处理相同起点，否则覆盖和去重会错。区间端点可能溢出时用更大整数。闭区间和半开区间决定是否用小于等于。\\
\end{tabular}
\end{principlebox}

\begin{principlebox}{题型分流：区间、排序与合并}
\textbf{合并和插入。}
判断条件：区间按起点排序后，重叠区间连续出现。处理方式：维护结果尾区间或新区间边界。代表案例：56 合并区间、57 插入区间。风险点：闭区间和半开区间端点相接处理不同。

\textbf{覆盖和删除。}
判断条件：判断一个区间是否被之前更大的右端覆盖。处理方式：同起点终点降序，扫描维护最大右端。代表案例：1288 删除被覆盖区间、435 无重叠区间。风险点：同起点排序方向是关键。

\textbf{调度和最少资源。}
判断条件：多个区间竞争时间资源，问最多保留或最少会议室。处理方式：按终点贪心或按起点扫描堆维护当前结束时间。代表案例：252 会议室、253 会议室 II、452 引爆气球。风险点：端点相等是否冲突由题意决定。

\textbf{动态区间。}
判断条件：区间逐个插入，必须立即判断冲突或维护覆盖。处理方式：有序表查前驱后继，或扫描线维护事件。代表案例：729 我的日程安排表 I、731 我的日程安排表 II。风险点：哈希表无法回答前驱后继。

\end{principlebox}

本节案例覆盖 LeetCode 56 合并区间、LeetCode 57 插入区间、LeetCode 252 会议室、LeetCode 253 会议室 II、LeetCode 352 将数据流变为多个不相交区间、LeetCode 436 寻找右区间、LeetCode 729 我的日程安排表 I、LeetCode 731 我的日程安排表 II、LeetCode 986 区间列表的交集、LeetCode 1094 拼车、LeetCode 1109 航班预订统计、LeetCode 1229 安排会议日程、LeetCode 1288 删除被覆盖区间。阅读时不要按题号背答案，而要先判断题面落在哪个分流场景：查询目标是什么，输入约束给了什么单调性或等价关系，暴力慢在哪里，优化结构保存了什么状态。

\subsection{代表案例与单点分析}

\subsubsection{LeetCode 56 合并区间}

\textbf{题目模型。} 排序后可能与当前区间相交的只会是结果尾部。

\textbf{推导重点。} 按起点排序，扫描时维护结果最后区间的右端点。

\textbf{边界。} 端点相接、完全覆盖、无重叠、空列表、输入未排序。

\textbf{变体追问。} 若区间是半开区间，端点相等不一定合并。

\subsubsection{LeetCode 57 插入区间}

\textbf{题目模型。} 新区间只会影响与它相交的一段连续区间。

\textbf{推导重点。} 先追加左侧完全不相交区间，再合并重叠段，最后追加右侧。

\textbf{边界。} 新区间在最前、最后、覆盖全部、刚好相邻、原列表为空。

\textbf{变体追问。} 若输入不是有序且不重叠，必须先排序或换模型。

\subsubsection{LeetCode 252 会议室}

\textbf{题目模型。} 每个会议占用一个时间区间，任意重叠都表示不能全部参加。

\textbf{推导重点。} 按开始时间排序，只需检查相邻会议是否冲突。

\textbf{边界。} 端点相等是否冲突、空列表、单会议、完全重叠。

\textbf{变体追问。} 若问最少会议室，需要维护同时进行的会议数量。

\subsubsection{LeetCode 253 会议室 II}

\textbf{题目模型。} 同时进行的会议数量决定最少房间数。

\textbf{推导重点。} 按开始时间扫描，用小顶堆维护当前会议的最早结束时间。

\textbf{边界。} 会议端点相接、多个同一时间开始、空列表、堆顶释放条件。

\textbf{变体追问。} 也可以把开始和结束分别排序，用双指针统计并发数。

\subsubsection{LeetCode 352 将数据流变为多个不相交区间}

\textbf{题目模型。} 数字逐个到达，需要动态维护已覆盖的有序不相交区间。

\textbf{推导重点。} 用有序表找到新值附近的前驱和后继，判断是否连接、合并或新建区间。

\textbf{边界。} 重复插入、连接左右两个区间、插入最小最大值、相邻端点。

\textbf{变体追问。} 这类动态区间不能用普通数组扫描长期维护。

\subsubsection{LeetCode 436 寻找右区间}

\textbf{题目模型。} 每个区间需要寻找起点不小于当前终点的最小起点区间。

\textbf{推导重点。} 把所有起点和原下标排序，对每个终点做 lower bound。

\textbf{边界。} 不存在右区间、起点相同、负端点、单区间。

\textbf{变体追问。} 如果区间动态插入，则需要有序表而不是一次排序数组。

\subsubsection{LeetCode 729 我的日程安排表 I}

\textbf{题目模型。} 每次插入新区间时，需要判断它是否和已有区间重叠。

\textbf{推导重点。} 有序表按起点存区间，只检查新区间的前驱和后继。

\textbf{边界。} 端点相接、插入最前最后、完全覆盖已有区间、动态多次插入。

\textbf{变体追问。} 哈希表无法回答前驱后继，动态区间要用有序结构。

\subsubsection{LeetCode 731 我的日程安排表 II}

\textbf{题目模型。} 允许双重预订但不允许三重预订，插入会改变区间重叠层数。

\textbf{推导重点。} 保存已有预订和已经双重预订的区间；新预订若碰到双重区间则失败，否则记录新产生的重叠。

\textbf{边界。} 端点相接、多个局部重叠、完全覆盖、回滚失败插入。

\textbf{变体追问。} 扫描线差分也能做，但要处理失败时恢复状态。

\subsubsection{LeetCode 986 区间列表的交集}

\textbf{题目模型。} 两个有序且不重叠的区间列表，交集只可能来自当前两个区间。

\textbf{推导重点。} 每次比较两个当前区间的重叠部分，然后推进右端较小的区间。

\textbf{边界。} 空列表、端点相接、一个区间覆盖多个区间、无交集。

\textbf{变体追问。} 这是区间双指针，不需要把两个列表合并排序。

\subsubsection{LeetCode 1094 拼车}

\textbf{题目模型。} 乘客上下车事件构成沿路线位置变化的载客量。

\textbf{推导重点。} 用差分数组或扫描线记录每个位置人数变化，累计人数不能超过容量。

\textbf{边界。} 同站上下车、容量刚好、站点范围、输入未排序。

\textbf{变体追问。} 航班预订统计也是差分思想，只是只需要最终每段累加结果。

\subsubsection{LeetCode 1109 航班预订统计}

\textbf{题目模型。} 每条预订给一个连续航班区间增加座位数。

\textbf{推导重点。} 差分数组在区间起点加人数，终点后一位减人数，最后前缀还原每航班总数。

\textbf{边界。} 区间到最后一个航班、单点区间、多条重叠预订、下标从一开始。

\textbf{变体追问。} 这是静态区间加法，不需要线段树。

\subsubsection{LeetCode 1229 安排会议日程}

\textbf{题目模型。} 两个人的可用时间段都有序，目标是找最早长度足够的交集。

\textbf{推导重点。} 双指针比较当前两个区间交集，若长度不足就推进结束更早的一侧。

\textbf{边界。} 无可行时间、刚好等于 duration、区间端点相接、输入需要先排序。

\textbf{变体追问。} 它和区间列表交集相同，但额外带最早可行和长度约束。

\subsubsection{LeetCode 1288 删除被覆盖区间}

\textbf{题目模型。} 被覆盖要求另一区间起点不晚且终点不早。

\textbf{推导重点。} 按起点升序、终点降序排序，扫描维护最大右端。

\textbf{边界。} 同起点区间、完全覆盖、无覆盖、端点相等。

\textbf{变体追问。} 终点降序是关键，否则同起点小区间会先出现造成误判。

